\section{链接变量的解析涂抹}
\label{sec:analytic}

一种在整个有限复平面上处处解析、因而可微的链接变量涂抹方法可以定义如下。
令 $C_\mu(x)$ 表示自格点 $x$ 出发、终止于相邻格点
$x\!+\!\hat{\mu}$ 的各垂直 staple 组合的带权和：
\begin{eqnarray}
 C_\mu(x)&=&\sum_{\nu\neq \mu}\rho_{\mu\nu}\biggl(
 U_\nu(x) U_\mu(x\!+\!\hat{\nu}) U_\nu^\dagger(x\!+\!\hat{\mu})\nonumber\\
&&+ U^\dagger_\nu(x\!-\!\hat{\nu}) U_\mu(x\!-\!\hat{\nu})
  U_\nu(x\!-\!\hat{\nu}\!+\!\hat{\mu})
\biggr), \label{eq:Cdef}
\end{eqnarray}
其中 $\hat{\mu},\hat{\nu}$ 分别是方向 $\mu,\nu$ 上长度为一个格距的矢量
（基本格点平移矢量）。权重 $\rho_{\mu\nu}$ 为可调实参数。然后，在 $SU(N)$
中由下式定义矩阵 $Q_\mu(x)$：
\begin{eqnarray}
Q_\mu(x) &=&
  \frac{i}{2}\Bigl(\!\Omega^\dagger_\mu(x)\!-\!\Omega_\mu(x)\!\Bigr)
   \!-\!\frac{i}{2N}\Tr\Bigl(\!\Omega^\dagger_\mu(x)
  \!-\!\Omega_\mu(x)\!\Bigr),\nonumber\\
\Omega_\mu(x) &=& C_\mu(x)\ U_\mu^\dagger(x),
   \quad\mbox{（对 $\mu$ 不求和）} \label{eq:Qdef}
\end{eqnarray}
该矩阵厄米且无迹，因而 $e^{i Q_\mu(x)}$ 是 $SU(N)$ 的元素。
我们利用这一点定义一种迭代的解析链接涂抹算法：把第 $n$ 步的链接
$U_\mu^{(n)}(x)$ 映射为第 $n+1$ 步的链接 $U_\mu^{(n+1)}(x)$，
\begin{equation}
U^{(n+1)}_\mu(x)=\exp\Bigl(i Q_\mu^{(n)}(x)\Bigr)\ U_\mu^{(n)}(x).
\label{eq:Ustout}
\end{equation}
$e^{iQ_\mu(x)}$ 属于 $SU(N)$ 这一事实保证 $U^{(n+1)}_\mu(x)$ 也属于
$SU(N)$，从而免去投影回规范群的必要。这一 fuzz 步骤可迭代 $n_\rho$ 次，
最终产生我们称之为 {\em stout 链接}\cite{pub} 的链接变量，记作
$\tilde{U}_\mu(x)$：
\begin{equation}
 U\rightarrow U^{(1)}\rightarrow U^{(2)}\rightarrow\cdots\rightarrow
  U^{(n_\rho)}\equiv \tilde{U}.
\end{equation}

staple 权重的一个常见选择是
\begin{equation}
   \rho_{jk}=\rho,\qquad \rho_{4\mu}=\rho_{\mu 4}=0,
\end{equation}
它给出只涂抹空间链接的三维方案。另一个常见选择是各向同性的四维方案，其中
所有权重取相同值 $\rho_{\mu\nu}=\rho$。注意，文献~\cite{luscher} 曾把这样
的方案用作场变换，以从最一般的 $O(a^2)$ on-shell 改进规范作用量中消去某个
相互作用项。

不难证明，只要恰当选取权重 $\rho_{\mu\nu}$，stout 链接就具有与原始链接
变量完全相同的对称变换性质。在任意局域规范变换 $G(x)$ 下，链接变量按
$U_\mu(x)\rightarrow G(x)\ U_\mu(x)\ G^\dagger(x\!+\!\hat{\mu})$ 变换，
其中 $G(x)$ 为 $SU(N)$ 矩阵。于是在这样的规范变换下，
\begin{eqnarray}
Q_\mu(x)&\rightarrow& G(x)\ Q_\mu(x)\ G^\dagger(x),\\
e^{iQ_\mu(x)}&\rightarrow& G(x)\ e^{iQ_\mu(x)}\ G^\dagger(x),\\
U^{(1)}_\mu(x)&\rightarrow& G(x)\ e^{iQ_\mu(x)} G^\dagger(x)
 G(x)U_\mu(x)\ G^\dagger(x\!+\!\hat{\mu})\nonumber\\
 &=& G(x)\ U^{(1)}_\mu(x)\ G^\dagger(x\!+\!\hat{\mu}),
\end{eqnarray}
正如所要求的。还可以容易地证明：若权重 $\rho_{\mu\nu}$
尊重格点的转动与反射对称性，则 stout 链接在所有转动以及含该链接的平面内
所有反射下都满足所需的变换性质。最后，考虑在垂直于链接 $U_\mu(x)$ 并通过
其中点 $x\!+\!\frac{1}{2}\hat{\mu}$ 的平面内的反射。在这种操作下，链接
$U_\mu(x)$ 按
$U_\mu(x)\rightarrow U_{-\mu}(x\!+\!\hat{\mu})=U_\mu^\dagger(x)$ 变换；
且当 $\nu\neq \mu$ 时，$U_\nu(x)\leftrightarrow U_\nu(x\!+\!\hat{\mu})$、
$U_\nu^\dagger(x\!+\!\hat{\mu})\leftrightarrow U_\nu^\dagger(x)$。
设 $\rho_{\mu\nu}$ 为实数，则得 $C_\mu(x)\rightarrow C_\mu^\dagger(x)$，
于是
\begin{eqnarray}
 \Omega_\mu(x) &\rightarrow & C_\mu^\dagger(x)U_\mu(x)=U_\mu^\dagger(x)
 \Omega_\mu^\dagger(x) U_\mu(x),\\
 Q_\mu(x)&\rightarrow & -U_\mu^\dagger(x)Q_\mu(x)U_\mu(x),\\
 U^{(1)}_\mu(x)&\rightarrow & U_\mu^\dagger(x) e^{-iQ_\mu(x)}
   = U_\mu^{(1)\dagger}(x),
\end{eqnarray}
正如所要求的。既然 $U^{(1)}_\mu(x)$ 在一切对称操作下的变换方式与
$U_\mu(x)$ 相同，便可得出结论：stout 链接 $\tilde{U}_\mu(x)$ 具有与原始
链接变量完全相同的对称变换性质。

\begin{figure}[tb]
\includegraphics[width=0.98\textwidth]{images/paths}
\caption[figF]{新链接变量 $U^{(1)}$ 按 $\rho_{\mu\nu}$ 展开到一阶时用原
链接路径表出的形式。每个闭合回路都带有一个迹，在 $SU(N)$ 中含因子
$1/N$。
\label{fig:paths}}
\end{figure}

由于指数函数具有收敛半径为无穷大的幂级数展开，每条 stout 链接都可视为
一个极其庞大而复杂的路径之和。当 $\rho_{\mu\nu}$ 较小时，把链接变量
$U^{(1)}$ 按 $\rho_{\mu\nu}$ 展开到一阶所得的路径示于
图~\ref{fig:paths}。注意，定义为
\begin{equation}
 U^{(1)}_\mu(x) = {\cal P}_{SU(3)}\ \biggl\{
 U_\mu(x) + C_\mu(x) \biggr\},
\label{eq:Ufuzz}
\end{equation}
的标准涂抹方法（其中 ${\cal P}_{SU(3)}$ 表示到 $SU(3)$ 的投影）在
$\rho_{\mu\nu}$ 的一阶处给出相同的路径和。

还有几点值得注意。首先，把权重 $\rho_{\mu\nu}$ 取为纯虚数的替代涂抹方案
并不能减小与理论高频模式的耦合。其次，不可能通过把作用量改用 stout 链接
表出来消除 Wilson 规范作用量的领头 $O(a^2)$ 离散误差；能做到的最佳程度是
通过减小蝌蚪图贡献来削弱 $O(\alpha_s a^2)$ 误差。第三，在 Symanzik 改进
作用量中使用 stout 链接，或许可以免除在确定作用量耦合系数时的蝌蚪改进。
最后，式~(\ref{eq:Cdef}) 给出的 $C_\mu(x)$ 定义并非唯一。按照上文概述的
方式，也可以用从格点 $x$ 到 $x\!+\!\hat{\mu}$ 的其他路径之和构造解析涂抹
方案，例如胖链接涂抹\cite{fat1} 与 HYP 涂抹\cite{hyp} 中所用的那些路径。
